\chapter{Single-Core Microarchitecture (\texttt{spCORE})}
\label{chap:core_microarchitecture}

\section{Structural Overview of the Compute Core}
Each of the ten identical graphics processing units (\textbf{\texttt{spCORE}}) instantiated in the top-level entity is designed as a modular, decoupled processing engine. As shown in Figure~\ref{fig:single_core}, \texttt{spCORE} is internally divided into two primary sub-blocks:
\begin{enumerate}
    \item \textbf{\texttt{spPIPE}} (\textit{Instruction Pipeline, Fetch, Decode \& Registers}): Manages the core command interface, decodes 64-bit spISA words, latches vertex and color parameters into dedicated registers, and orchestrates the core execution state machine.
    \item \textbf{\texttt{spEXEC}} (\textit{Execution Unit, Dedicated Accelerators \& Tile Boundary Protection}): Houses the parallel hardware rasterization engines, arbitrates the generated pixel streams, handles swap handshaking, and enforces hard spatial clipping boundaries for the tile.
\end{enumerate}

\begin{figure}[htbp]
\centering
\resizebox{0.95\textwidth}{!}{
    \begin{tikzpicture}[
    >=Stealth,
    node distance=1.0cm and 1.2cm,
    every text node part/.style={align=center},
    box/.style={draw=black!80, thick, rounded corners=3pt, fill=white, inner sep=5pt, font=\small},
    pipebox/.style={box, fill=blue!8, draw=blue!60},
    execbox/.style={box, fill=emerald!8, draw=emerald!60},
    accbox/.style={box, fill=amber!15, draw=amber!80, minimum width=2.6cm, minimum height=0.75cm, font=\scriptsize\bfseries},
    clipbox/.style={box, fill=red!10, draw=red!70, font=\small\bfseries},
    bus/.style={draw=black!70, line width=1.3pt, ->},
    sig/.style={draw=blue!70!black, thick, ->},
    ctl/.style={draw=purple!80!black, dashed, thick, ->},
    subframe/.style={draw=gray!60, dashed, rounded corners=6pt, inner sep=8pt}
]

% ==================== spCORE INPUTS ====================
\node[box, fill=gray!20] (in_fifo) {
    \textbf{From Scheduler FIFO}\\
    \scriptsize\texttt{instr\_word[63:0]}\\
    \scriptsize\texttt{instr\_valid}
};

% ==================== spPIPE UNIT ====================
\begin{scope}[local bounding box=pipe_unit]
    \node[pipebox, right=1.2cm of in_fifo, minimum width=3.4cm] (fsm_fetch) {
        \textbf{Pipeline FSM \& Fetch}\\
        \scriptsize States: \texttt{normal}, \texttt{drawing}, \texttt{halt}\\
        \scriptsize Handshake: \texttt{instr\_req}
    };

    \node[pipebox, below=0.8cm of fsm_fetch, minimum width=3.4cm] (decoder) {
        \textbf{Opcode Decoder}\\
        \scriptsize Decodes \texttt{OP[3:0]}\\
        \scriptsize Generates \texttt{acc\_enable\_vec[5:0]}
    };

    \node[pipebox, below=0.8cm of decoder, minimum width=3.4cm] (param_regs) {
        \textbf{Parameter Registers}\\
        \scriptsize $X_1, Y_1, X_2, Y_2, X_3, Y_3$ (9-bit)\\
        \scriptsize Depth $Z$ (4-bit), Color (15-bit RGB)
    };

    \draw[bus] (in_fifo) -- (fsm_fetch);
    \draw[sig] (fsm_fetch) -- (decoder);
    \draw[bus] (decoder) -- (param_regs);
\end{scope}
\node[subframe, fit=(pipe_unit), label={[anchor=north west, font=\bfseries\color{blue!80!black}]above left:\texttt{spPIPE} (Instruction Fetch, Decode \& Registers)}] {};

% ==================== ACCELERATORS (in spEXEC) ====================
\begin{scope}[local bounding box=acc_array]
    \node[accbox, right=2.0cm of fsm_fetch.north, anchor=north west] (acc_pixel) {
        DRAWPIXEL Engine (0)\\
        \tiny 1-Cycle Direct Latch
    };
    \node[accbox, below=0.25cm of acc_pixel] (acc_line) {
        LINE\_ACC (1)\\
        \tiny Bresenham Line ($\Delta X, \Delta Y$, err)
    };
    \node[accbox, below=0.25cm of acc_line] (acc_tri) {
        EDGE\_FILL\_TOP\_V2 (3)\\
        \tiny Triangle 2D Edge Cross-Products
    };
    \node[accbox, below=0.25cm of acc_tri] (acc_circ) {
        CIRCLE\_ACC (4)\\
        \tiny Midpoint Circle (8 Octants)
    };
    \node[accbox, below=0.25cm of acc_circ] (acc_fcirc) {
        F\_CIRCLE\_ACC (5)\\
        \tiny Filled Circle (Bresenham + Spans)
    };
\end{scope}

% Routing from spPIPE to Accelerators
\draw[ctl] (decoder.east) -- ++(0.6,0) |- node[pos=0.25, right, font=\tiny] {\texttt{acc\_enable\_vec}} (acc_pixel.west);
\draw[ctl] (decoder.east) -- ++(0.6,0) |- (acc_line.west);
\draw[ctl] (decoder.east) -- ++(0.6,0) |- (acc_tri.west);
\draw[ctl] (decoder.east) -- ++(0.6,0) |- (acc_circ.west);
\draw[ctl] (decoder.east) -- ++(0.6,0) |- (acc_fcirc.west);

\draw[bus] (param_regs.east) -- ++(0.9,0) |- node[pos=0.25, right, font=\tiny] {Coords, Color} (acc_pixel.195);
\draw[bus] (param_regs.east) -- ++(0.9,0) |- (acc_line.195);
\draw[bus] (param_regs.east) -- ++(0.9,0) |- (acc_tri.195);
\draw[bus] (param_regs.east) -- ++(0.9,0) |- (acc_circ.195);
\draw[bus] (param_regs.east) -- ++(0.9,0) |- (acc_fcirc.195);

% ==================== EXECUTION & ARBITRATION ====================
\begin{scope}[local bounding box=exec_tail]
    % Pixel Multiplexer
    \node[box, fill=teal!15, draw=teal!80, right=1.2cm of acc_tri, minimum height=3.5cm, minimum width=1.6cm] (mux_pixel) {
        \textbf{Candidate}\\
        \textbf{Pixel Mux}\\
        \tiny Selects active\\
        \tiny accelerator\\
        \tiny based on\\
        \tiny\texttt{acc\_busy\_vec}
    };

    \draw[bus] (acc_pixel.east) -- (acc_pixel.east -| mux_pixel.west);
    \draw[bus] (acc_line.east) -- (acc_line.east -| mux_pixel.west);
    \draw[bus] (acc_tri.east) -- (mux_pixel.west);
    \draw[bus] (acc_circ.east) -- (acc_circ.east -| mux_pixel.west);
    \draw[bus] (acc_fcirc.east) -- (acc_fcirc.east -| mux_pixel.west);

    % Completion OR Gate
    \node[box, fill=orange!15, draw=orange!80, below=0.6cm of mux_pixel, minimum width=2.2cm] (or_finish) {
        \textbf{OR Aggregator}\\
        \tiny\texttt{finish\_exec <= OR(acc\_finish)}
    };
    \draw[sig] (or_finish.west) to[out=180, in=-90] node[below, font=\tiny] {\texttt{finish\_exec}} (fsm_fetch.south);

    % Hardware Tile Clipping Unit
    \node[clipbox, right=1.2cm of mux_pixel, minimum width=3.4cm, minimum height=1.6cm] (clip_unit) {
        \textbf{Tile Clipping Unit}\\
        \scriptsize $Y_{\min} = \text{tile\_index} \times 24$\\
        \scriptsize $Y_{\max} = Y_{\min} + 23$\\
        \scriptsize If $Y \in [Y_{\min}, Y_{\max}]$: \texttt{valid = 1}\\
        \scriptsize Else: \texttt{pixel\_valid\_o = 0}
    };
    \draw[bus] (mux_pixel) -- node[above, font=\scriptsize] {$X, Y, C, Z$} node[below, font=\tiny] {\texttt{pixel\_valid\_int}} (clip_unit);

    \node[box, fill=purple!15, draw=purple!80, right=1.2cm of clip_unit] (out_tile) {
        \textbf{To Tile VRAM $i$}\\
        \scriptsize\texttt{pixel\_x\_o[8:0]}\\
        \scriptsize\texttt{pixel\_y\_o[8:0]}\\
        \scriptsize\texttt{pixel\_color\_o[14:0]}\\
        \scriptsize\texttt{pixel\_z\_o[3:0]}\\
        \scriptsize\texttt{pixel\_valid\_o}
    };
    \draw[bus] (clip_unit) -- (out_tile);
\end{scope}

\node[subframe, fit=(acc_array) (exec_tail), label={[anchor=north west, font=\bfseries\color{emerald!80!black}]above left:\texttt{spEXEC} (Execution Unit, Dedicated Accelerators \& Tile Boundary Protection)}] {};

\end{tikzpicture}

}
\caption{Microarchitecture of the Single Graphics Core (\texttt{spCORE}).}
\label{fig:single_core}
\end{figure}

\section{Pipeline Control and Decode Unit (\texttt{spPIPE.vhd})}
The \texttt{spPIPE} block is responsible for synchronizing instruction ingestion with accelerator availability.

\subsection{Pipeline Finite State Machine}
The pipeline controller is governed by an FSM with three operational states:
\begin{itemize}
    \item \texttt{normal}: The core asserts \texttt{instr\_req <= '1'} toward its dedicated scheduler FIFO. When \texttt{instr\_valid} is sampled high, the 64-bit instruction word is latched. The opcode decoder evaluates bits \texttt{[3:0]}:
    \begin{itemize}
        \item If \texttt{NOP}, parameters are cleared and the core remains in \texttt{normal}.
        \item If \texttt{SETCOLOR}, the 15-bit color field is stored in the internal color register \texttt{color\_int} and the state remains \texttt{normal}.
        \item If a drawing primitive or \texttt{SWAP\_BUFFERS} is detected, the corresponding accelerator enable line in \texttt{acc\_enable\_vec[5:0]} is pulsed, the busy mask \texttt{acc\_busy\_vec} is set, and the FSM transitions to \texttt{drawing}.
    \end{itemize}
    \item \texttt{drawing}: In this state, \texttt{instr\_req} is held low to stall further instruction fetching. The core remains in \texttt{drawing} until the execution unit signals completion via \texttt{finish\_exec = '1'}, at which point the state transitions back to \texttt{normal}.
    \item \texttt{halt}: An asynchronous stall state entered whenever the external signal \texttt{core\_halt} is asserted.
\end{itemize}

\subsection{Parameter Registration}
During decode, the packed bitfields of the instruction word are extracted and registered into dedicated 9-bit coordinate buses: $(X_1, Y_1)$, $(X_2, Y_2)$, $(X_3, Y_3)$, depth $Z$, and the active 15-bit RGB color.

\section{Execution Unit and Tile Boundary Protection (\texttt{spEXEC.vhd})}
The \texttt{spEXEC} module interfaces the control signals from \texttt{spPIPE} to the physical rasterization engines.

\subsection{Candidate Pixel Multiplexing and Completion Aggregation}
When a drawing primitive is active, its corresponding hardware accelerator outputs a continuous stream of candidate pixels $(X, Y, \text{Color})$ qualified by an individual \texttt{pixel\_valid} flag. The multiplexer process \texttt{pixel\_out\_proc} routes the active accelerator's outputs to the internal bus based on \texttt{acc\_busy\_vec}.

Simultaneously, the completion signals from all six accelerators and the swap controller are logically OR-ed to generate \texttt{finish\_exec}:
\begin{equation}
\texttt{finish\_exec} = \bigvee_{k=0}^{5} \texttt{acc\_finish\_vec}(k) \lor \texttt{finish\_swap}
\end{equation}
This single-cycle pulse informs \texttt{spPIPE} that the primitive has been completely rasterized, allowing the core to fetch the next instruction from its FIFO.

\subsection{Hardware Tile Clipping Unit}
While the upstream scheduler restricts instruction dispatch to overlapping tiles, primitives that cross multiple tiles will naturally produce pixels outside the vertical bounds of any single tile. 

To prevent out-of-bounds pixels from contaminating adjacent tile memory blocks or causing address wrap-around, \texttt{spEXEC} implements a dedicated, zero-latency \textbf{Tile Clipping Process} (\texttt{pixel\_clip\_proc}):
\begin{lstlisting}[language=VHDL, caption={Tile Boundary Enforcement in \texttt{spEXEC.vhd}.}, label={lst:tile_clipping}]
pixel_clip_proc : process(pixel_x_int, pixel_y_int, pixel_color_int, z_i, pixel_valid_int, tile_index)
    variable x_coord, y_coord : integer;
    variable y_min, y_max     : integer;
begin
    x_coord := to_integer(unsigned(pixel_x_int));
    y_coord := to_integer(unsigned(pixel_y_int));
    
    y_min := to_integer(unsigned(tile_index)) * 24;
    y_max := y_min + 23;
    
    pixel_x_o     <= pixel_x_int;
    pixel_y_o     <= pixel_y_int;
    pixel_color_o <= pixel_color_int;
    pixel_z_o     <= z_i;
    
    if pixel_valid_int = '1' then
        if (x_coord >= 0 and x_coord < 320) and (y_coord >= y_min and y_coord <= y_max) then
            pixel_valid_o <= '1';
        else
            pixel_valid_o <= '0'; -- Zero-overhead hardware suppression
        end if;
    else
        pixel_valid_o <= '0';
    end if;
end process pixel_clip_proc;
\end{lstlisting}

If a candidate pixel's $Y$ coordinate satisfies $Y \in [i \times 24, i \times 24 + 23]$, it is written into VRAM. If it falls outside this 24-line boundary, \texttt{pixel\_valid\_o} is driven to \texttt{'0'}, silently discarding the pixel without interrupting the accelerator pipeline.
